\chapter*{Tóm tắt}
Nội dung của đồ án hướng đến việc nghiên cứu, thiết kế và hiện thực một IoT gateway dạng mô-đun có khả năng khả cấu hình linh hoạt, nhằm giải quyết các hạn chế của các giải pháp gateway truyền thống trong các hệ thống IoT hiện nay. Trong bối cảnh mỗi kịch bản ứng dụng IoT đòi hỏi tập hợp giao thức và chuẩn giao tiếp khác nhau, việc phát triển các gateway tích hợp sẵn nhiều chức năng không cần thiết dẫn đến lãng phí tài nguyên và làm tăng chi phí triển khai.

Đề tài tập trung tìm hiểu các khái niệm nền tảng liên quan đến IoT gateway, làm rõ vai trò của gateway trong hạ tầng IoT, đồng thời phân tích các công nghệ, giao thức và chuẩn giao tiếp phổ biến đang được sử dụng. Trên cơ sở đó, một kiến trúc gateway định hướng mô-đun được đề xuất, trong đó phần cứng được thiết kế theo dạng lắp ghép nhằm cho phép lựa chọn các module phù hợp với từng kịch bản ứng dụng, trong khi phần mềm hỗ trợ quy trình cấu hình tối giản và thân thiện với người dùng.

Trong giai đoạn hiện thực, đề tài tiến hành xây dựng nguyên mẫu (prototype) của hệ thống nhằm kiểm chứng tính khả thi của kiến trúc đề xuất. Các thử nghiệm ban đầu cho thấy hệ thống hoạt động ổn định, đáp ứng được các yêu cầu chức năng cơ bản và thể hiện tiềm năng tối ưu hóa chi phí triển khai so với các giải pháp tích hợp truyền thống. Kết quả đạt được tạo nền tảng cho việc tiếp tục hoàn thiện và mở rộng hệ thống trong các giai đoạn phát triển sản phẩm tiếp theo, hướng tới một IoT gateway có khả năng ứng dụng trong thực tế.